\clearpage
\chapter*{Conclusiones}
\phantomsection
\addcontentsline{toc}{chapter}{Conclusiones}

A lo largo de la ejecución las evidencias se ha completado el ciclo de especificación y validación de los requisitos preliminares del proyecto.

Se ha demostrado la capacidad de transmitir las necesidades del negocio, capturadas en \textbf{Historias de Usuario}, a un lenguaje de ingeniería formal a través de \textbf{Requisitos Funcionales (RF) y No Funcionales (RNF)}, enmarcado en el estándar \textbf{IEEE 830}. La selección y configuración de \textbf{Notion} como herramienta de gestión ha probado ser una decisión estratégica acertada, permitiendo una trazabilidad robusta que herramientas más simples no ofrecen.

Finalmente, la creación de \textbf{prototipos y casos de prueba} ha permitido ``aterrizar'' los conceptos abstractos, validando que las interfaces propuestas son lógicamente coherentes y funcionales. Aunque estos artefactos son de baja fidelidad, constituyen una prueba de concepto esencial que mitiga significativamente el riesgo de desarrollar funcionalidades que no se alineen con la visión del producto.
